\section{One Construction for Every Notation}\label{sec:general-framework}

One construction turns a process model into a signature of generator cospans, and it is the
same for every notation in Sections~\ref{sec:petri-nets}--\ref{sec:bpmn}. Pick an activity,
follow the edges around it through the mediators, and stop at the first activities met on each
side, recording the object types and edge weights seen on the way. Each way of joining what the
activity receives to what it emits is one generator. Composing generators freely gives a
symmetric monoidal category, each of its morphisms is a partial order of activity occurrences,
and its traces are the orderings of those occurrences. For any choice of morphisms made from the
signature alone, models with equal signatures have equal languages. A badly specified model appears as composites that cannot be satisfied.

\subsection{The AND/OR graph}
\label{subsec:firing-graphs}

Every notation is first mapped to one kind of graph. Its nodes are activities, which a trace
records, and mediators, which it does not. A mediator sends flow along all its
branches at once (AND) or along exactly one of them (XOR). Edges carry the typing and weighting
that the notation records.

\begin{definition}[AND/OR graph]
  \label{def:lm-graph}
  An AND/OR graph over a finite set $\mathcal{O}$ of object types is a tuple
  $G=(\mathcal{A},R,E,\ell,\nu,\mathrm{type},\mathrm{wt},\rho)$ in which $(\mathcal{A}\sqcup R,E)$
  is a finite directed graph whose nodes are the activities $\mathcal{A}$ and the mediators $R$,
  $\ell:\mathcal{A}\to\mathcal{L}$ labels the activities, $\nu:R\to\{\mathrm{AND},\mathrm{XOR}\}$
  gives each mediator its kind, $\mathrm{type}:E\to\mathcal{O}$ and
  $\mathrm{wt}:E\to\mathbb{N}_{\ge 1}\cup\{\mathrm{var}\}$ give each edge a type and a weight, and
  $\rho(a)$ is a set of linear equalities and inequalities among the weights of the edges at the
  activity $a$. The weight $\mathrm{var}$ marks an edge that carries zero or more objects, and
  $\rho(a)$ is empty unless the notation relates the weights. Silent transitions are mediators,
  and a notation with no types takes $\mathcal{O}=\{\ast\}$.
\end{definition}

Object-centric Petri nets, object-centric causal nets, process trees and BPMN are instances.
The work specific to a notation is the correctness of its map to AND/OR graphs, which each
notation's section establishes. Everything below is the same for all of them.

\subsection{Reach families, contexts and generators}
\label{subsec:generators}

\begin{definition}[Reach families and contexts]
  \label{def:contexts-general}
  Let $a\in\mathcal{A}$. For a node $x$ reached from $a$ along a path $\pi$ of out-edges, let
  $w_\pi\subseteq\mathcal{O}$ be the set of types on the edges of $\pi$ and let $N(x)$ be the
  out-neighbours $y$ of $x$ with $R_a(y,\pi y)\ne\emptyset$. The reached family $R_a(x,\pi)$, a set of
  finite sets of typed wires, is
  \begin{gather*}
    R_a(x,\pi)=\twocolbreak
    \begin{cases}
      \{\{(a,w_\pi,x)\}\} & x\in\mathcal{A},\\
      \emptyset & x\text{ earlier on }\pi,\\
      \{\{(a,w_\pi,\top)\}\} & \text{no out-edges},\\
      \textstyle\bigcup_{y\in N(x)}R_a(y,\pi y) & \nu(x)=\mathrm{XOR},\\
      \{\textstyle\bigcup_{y\in N(x)}S_y : S_y\in R_a(y,\pi y)\} & \nu(x)=\mathrm{AND},
    \end{cases}
  \end{gather*}
  taking the cases in order, so the last four apply to mediators, with the AND case empty when
  $N(x)$ is. The forward family $\mathcal{F}_a$ joins the out-edges of
  $a$ as an AND mediator does, and $\mathcal{F}_a=\{\emptyset\}$ when no out-edge reaches anything.
  The backward family $\mathcal{B}_a$ is the same construction along in-edges, with typed wires
  $(x,w_\pi,a)$ and $(\bot,w_\pi,a)$. The contexts of $a$ are the pairs
  $\mathrm{Contexts}(a)=\mathcal{B}_a\times\mathcal{F}_a$.
\end{definition}

The walk visits each node at most once per path, so $\mathcal{F}_a$ and $\mathcal{B}_a$ are
finite on every finite graph, cycles included. Activities are identified by their labels, so
typed wires are compared by label.

\begin{definition}[Generator cospan]
  \label{def:generator-cospan-general}
  For $a\in\mathcal{A}$, each context $c=(P,S)\in\mathrm{Contexts}(a)$ gives the generator
  $g_{a,c}=(P\hookrightarrow P+S\hookleftarrow S,\ d_{a,c})$ with one hyperedge $e_a$, labelled
  $\lambda(e_a)=\ell(a)$, whose source is $P$ and whose target is $S$. Its decoration $d_{a,c}$
  also carries the constraint system $\Lambda_{a,c}$ of
  Definition~\ref{def:multiplicity-decoration}. Every context gives a generator and none is
  discarded. The signature of $G$ is
  $\Sigma=\{g_{a,c} : a\in\mathcal{A},\ c\in\mathrm{Contexts}(a)\}$.
\end{definition}

A typed wire is identified by its triple $(a,w,b)$. The set of edge weights met on its paths is
data of the typed wire and plays no part in its identity, and each generator's constraints use the weight of
the edge at its own end. The counts of a generator are fixed by its own legs only. Nothing
outside the generator ties its input counts to its output counts, so an activity may create or
consume objects of a type that appears on one side only.

\begin{definition}[Constraint system]
  \label{def:multiplicity-decoration}
  Each leg $p\in P\cup S$ of $g_{a,c}$ has a count $n_p\in\mathbb{N}$, the number of objects
  on $p$ in one occurrence of $a$. If the edge at $a$'s end of the path has weight $k$ then
  $n_p=k$, and if that edge is variable then $n_p\ge 0$. The equalities and inequalities
  $\rho(a)$ of Definition~\ref{def:lm-graph} apply to the counts of the legs whose edge at $a$'s
  end they name, as the shared keys and cardinality intervals of object-centric causal nets do. The constraint system $\Lambda_{a,c}$ is the conjunction of these
  constraints. Two generators are equal when they have the same label, the same typed wires and
  constraint systems with the same solutions.
\end{definition}

\subsection{The free category and its traces}
\label{subsec:free-category}

A leg with count $n$ is the $n$-fold tensor power of its typed wire, with $n$ ports. A port joins
an output of one occurrence to an input of another on the same typed wire, and nothing else
composes.

\begin{definition}[Free category of a signature]
  \label{def:free-category}
  For $g=g_{a,c}\in\Sigma$ and a count $k_p\in\mathbb{N}$ for each leg $p$, the count assignment of
  $g$ with these counts is the box
  \[
    g^{k}:\ \textstyle\bigotimes_{p\in P}p^{\otimes k_p}\ \to\ \bigotimes_{q\in S}q^{\otimes k_q}
  \]
  labelled $\ell(a)$, with $k_p$ ports on each leg $p$ in a fixed order of the legs.
  $\mathbf{F}(\Sigma)$ is the free symmetric monoidal category whose objects are the finite lists of typed wires and
  whose generating morphisms are all count assignments. A morphism $f$ is satisfiable when every count
  assignment $g^k$ in it satisfies $\Lambda_{a,c}$ with $n_p=k_p$.
\end{definition}

A morphism of $\mathbf{F}(\Sigma)$ is a string diagram built from count assignments, identities and
symmetries by $;$ and $\otimes$. A count assignment is an activity morphism of
Section~\ref{sec:hypergraphs} whose multiplicities on each typed wire are the counts $k_p$, so each
generator of $\mathbf{H}(\Sigma)$ gives one count assignment per choice of counts. $\mathbf{F}(\Sigma)$ is
the Frobenius-free part of $\mathbf{H}(\Sigma)$, the diagrams in which no port is copied, merged,
created or discarded.

A morphism belongs to $\mathbf{F}(\Sigma)$ whatever its boundary. An open input port stands
for an object that enters from outside, and therefore every initialisation of the model is some
morphism. A loop in the model appears as repeated composition, since a diagram has no cycles.

\begin{definition}[Trace semantics of a morphism]
  \label{def:trace-map}
  The occurrences of a morphism $f$ of $\mathbf{F}(\Sigma)$ are its count assignments. An occurrence $e$ precedes an occurrence
  $e'$ when a port of $f$ joins an output of $e$ to an input of $e'$, and the causal partial
  order $D_f$ is the reflexive and transitive closure of this relation. If $f$ is satisfiable then
  \begin{gather*}
    \Gamma(f)=\{\ell(e_1)\cdots\ell(e_m) :\twocolbreak
      e_1,\dots,e_m \text{ a linear extension of } D_f\},
  \end{gather*}
  where a silent label contributes the empty word, and otherwise $\Gamma(f)=\emptyset$.
\end{definition}

The relation $D_f$ is a partial order because diagrams are acyclic. A composite that contains
an unsatisfiable count assignment is itself unsatisfiable, and $\Gamma$ is empty on it.

\begin{exmp}[An AND split and join]
  \label{ex:and-fanout}
  Let $a$ lead through an AND mediator to $b$ and $c$, and let both lead through a second AND mediator to
  $d$, with one type and weight $1$ throughout. Write $(a,b)$ for $(a,\{\ast\},b)$, and write an
  activity's name for the count assignment of its generator with one port on each leg, and $I$ for the
  monoidal unit, the empty list of typed wires. The walk gives
  $\mathcal{F}_a=\{\{(a,b),(a,c)\}\}$ and $\mathcal{B}_d=\{\{(b,d),(c,d)\}\}$, and the count assignments include
  \[
    a:\ I\to(a,b)\otimes(a,c),\qquad b:\ (a,b)\to(b,d).
  \]
  With $c:(a,c)\to(c,d)$ and $d:(b,d)\otimes(c,d)\to I$, the composite
  $f=a\,;(b\otimes c);d$ has the order $a\le b\le d$ and $a\le c\le d$. Its two linear extensions
  give $\Gamma(f)=\{abcd,\ acbd\}$. The open composite $a\,;(b\otimes\mathrm{id}_{(a,c)})$ also
  belongs to $\mathbf{F}(\Sigma)$ and gives $\{ab\}$. Had the edge leaving $a$ carried weight $2$,
  the count assignment of $a$ with one port on each leg would be unsatisfiable and $\Gamma(f)=\emptyset$
  for this $f$, while the count assignment of $a$ with two ports on each leg composes.
\end{exmp}

\begin{lemma}[$\Gamma$ is well defined]
  \label{lem:gamma-well-defined}
  Equal morphisms of $\mathbf{F}(\Sigma)$ have equal images under $\Gamma$.
\end{lemma}
\begin{proof}
  A morphism is a class of terms under the symmetric monoidal equations, and $\Gamma$ depends on a
  term only through its occurrences, their labels and which output port reaches which input. Identities
  and symmetries have no occurrences and only relay ports. The associativity and unit laws regroup a
  term without adding occurrences or redirecting ports. Interchange, $(f\otimes g);(f'\otimes g')=
  (f;f')\otimes(g;g')$, joins the same outputs to the same inputs on both sides, and naturality
  of the symmetry moves a crossing past a box with every port ending where it did. The law
  $\sigma;\sigma=\mathrm{id}$ and the hexagons relay each port unchanged. Each equation therefore
  preserves the occurrences, satisfiability and the wiring, which records which output reaches
  which input, the boundary included. The wiring of $f;g$ and of $f\otimes g$ is fixed by
  those of $f$ and $g$, so the invariance extends to every term containing an equated part, and
  $D_f$ is determined by the wiring.
\end{proof}

\subsection{The model invariant}
\label{subsec:theorem}
\label{subsec:minimal-signatures}

A selection $\mathcal{S}$ assigns to each signature $\Sigma$ a set $\mathcal{S}(\Sigma)$ of morphisms of
$\mathbf{F}(\Sigma)$ by a rule that uses $\Sigma$ alone. Its language is
$L_{\mathcal{S}}(\Sigma)=\bigcup_{f\in \mathcal{S}(\Sigma)}\Gamma(f)$.

\begin{theorem}[Model invariant]
  \label{thm:canonical-presentation}
  For every selection $\mathcal{S}$, if $\Sigma_1=\Sigma_2$ then $L_{\mathcal{S}}(\Sigma_1)=L_{\mathcal{S}}(\Sigma_2)$.
\end{theorem}
\begin{proof}
  Equal signatures have the same typed wires and the same count assignments, with the same labels and the same
  satisfiable counts, so $\mathbf{F}(\Sigma_1)=\mathbf{F}(\Sigma_2)$. A rule that uses the
  signature alone selects the same morphisms from both, and by
  Lemma~\ref{lem:gamma-well-defined} the unions defining $L_{\mathcal{S}}$ agree.
\end{proof}

The theorem holds for every selection, and how well the language separates models depends on the
selection. The default selection takes the connected process diagrams of Section~\ref{sec:hypergraphs},
the morphisms with any boundary that have at least one count assignment and whose count assignments are
joined into one piece by their ports. It covers every
initialisation, including a pure cycle $a\to b\to c\to a$ started at whichever activity receives the first object, and it
keeps the order of occurrences, whereas over all of $\mathbf{F}(\Sigma)$ count assignments placed side by side
give every ordering of the labels. The converse of the theorem fails, since different signatures
can give equal languages.

Signatures can also be compared by inclusion. Write $\Sigma_1\subseteq\Sigma_2$ when every generator of
$\Sigma_1$ is equal to a generator of $\Sigma_2$ in the sense of
Definition~\ref{def:multiplicity-decoration}, with the same label, the same typed wires and
constraint systems with the same solutions.

\begin{proposition}[Inclusion of free categories]
  \label{prop:inclusion-functor}
  If $\Sigma_1\subseteq\Sigma_2$ then there is a strict symmetric monoidal functor
  $\mathrm{inc}\colon\mathbf{F}(\Sigma_1)\to\mathbf{F}(\Sigma_2)$ that is the identity on typed wires
  and on count assignments, and $\Gamma(\mathrm{inc}\,f)=\Gamma(f)$ for every morphism $f$.
\end{proposition}
\begin{proof}
  Every typed wire of $\Sigma_1$ is a typed wire of $\Sigma_2$, so every object of
  $\mathbf{F}(\Sigma_1)$ is an object of $\mathbf{F}(\Sigma_2)$. A count assignment $g^{k}$ of
  $\Sigma_1$ has the label, legs and counts of the count assignment $g^{k}$ of the equal generator of
  $\Sigma_2$, and it satisfies one constraint system exactly when it satisfies the other, since the
  two have the same solutions. Sending each generating morphism of $\mathbf{F}(\Sigma_1)$ to the same
  generating morphism of $\mathbf{F}(\Sigma_2)$ extends to a strict symmetric monoidal functor by
  the universal property of the free symmetric monoidal category. The functor sends a term built by
  $;$ and $\otimes$ from count assignments, identities and symmetries to the same term, so
  $\mathrm{inc}\,f$ has the occurrences of $f$ with the same labels and the same satisfiability, and
  each of its ports joins the output and input it joins in $f$. Hence $D_{\mathrm{inc}\,f}=D_f$ and
  $\Gamma(\mathrm{inc}\,f)=\Gamma(f)$. By Lemma~\ref{lem:gamma-well-defined} the equation holds for
  every term representing the morphism.
\end{proof}

A selection $\mathcal{S}$ is monotone when $\Sigma_1\subseteq\Sigma_2$ implies
$\mathrm{inc}(\mathcal{S}(\Sigma_1))\subseteq\mathcal{S}(\Sigma_2)$.

\begin{corollary}[Signature inclusion]
  \label{cor:signature-inclusion}
  For every monotone selection $\mathcal{S}$, if $\Sigma_1\subseteq\Sigma_2$ then
  $L_{\mathcal{S}}(\Sigma_1)\subseteq L_{\mathcal{S}}(\Sigma_2)$.
\end{corollary}
\begin{proof}
  A word of $L_{\mathcal{S}}(\Sigma_1)$ lies in $\Gamma(f)$ for some $f\in\mathcal{S}(\Sigma_1)$. By
  Proposition~\ref{prop:inclusion-functor} it lies in $\Gamma(\mathrm{inc}\,f)$, and
  $\mathrm{inc}\,f\in\mathcal{S}(\Sigma_2)$ because $\mathcal{S}$ is monotone.
\end{proof}

The connected process diagrams form a monotone selection. Having at least one count assignment
and having all count assignments joined into one piece by ports are properties of the occurrences
and the wiring, and $\mathrm{inc}$ preserves both. The corollary needs monotonicity, whereas
Theorem~\ref{thm:canonical-presentation} holds for every selection. The rule that selects the
morphisms in which every generator of $\Sigma$ occurs uses $\Sigma$ alone, yet a larger signature
can lose morphisms under it. The corollary gives no strictness, since a proper inclusion of
signatures may still give equal languages.

\subsection{Analysis of composites}
\label{subsec:analysis}

The analysis of a particular model selects and diagnoses composites of $\mathbf{F}(\Sigma)$, and
the \texttt{proc-posets} implementation performs it without changing $\Sigma$. A start generator
provides objects for one way the process can begin, a set of first activities that receive
objects together, and an end generator takes objects at the matching finish. These generators
form $\Sigma_\gamma$, every report states them as an assumption, and a selection that uses them
needs $\Sigma_\gamma$ equal as well as $\Sigma$ for Theorem~\ref{thm:canonical-presentation} to
apply. It is monotone when the start and end generators of the smaller signature are among those
of the larger, since a closed run stays closed under $\mathrm{inc}$, and
Corollary~\ref{cor:signature-inclusion} then applies. A closed run joins a start generator to an end generator with no open port. A livelock
family is a sequence of composites that another pass of a loop can always extend and that never
closes. A base run is a closed run from which no loop pass can be removed. Removal is
read on labels, so a closed run is a base run when no shorter closed run with the same start and
end generators differs from it, as label counts, by a sum of loop passes. Every closed run is then a
base run with loop passes inserted, and the closed runs fall into classes indexed by how many
passes are inserted where (Section~\ref{sec:validation}). Loops are classed by how one pass changes the object counts. A bounded loop leaves them
unchanged, a growing but closable loop adds objects that an end generator still absorbs, an
explosive loop produces objects that nothing can absorb, and a livelock loop is never left. Run
coverage lists the generators that lie on no run and the reason for each. When counts cannot
balance from the stated start, the report gives the least extra supply that would balance them.
Open typed wires that nothing joins, typed wires whose paths carry more than one type and models whose
activities form separate components are reported as well.

The \texttt{proc-posets} test suite checks the extraction on 26 hand-derived models against an
independent token game. On its bench, equal signatures never gave different languages under any
selection tried, and the connected process diagrams gave the fewest pairs of models with different
signatures and equal languages.
